\section{Finer row-sensitive character ranges}
\label{app:finer-character-bounds}

The uniform character theorem needs only the Schur relation, quadratic
reduction, bivariate and hook exactness, and one cubic witness. This
appendix retains the separate row-sensitive trace bounds and their
closed boundary, which can give better ranges for individual rows.

\subsection{Row-sensitive support of moment Koszul homology}

The first row of the row partition gives a second measure of how
early a tail relation can occur.  A relation on $t+1$ positive
inputs has many linear inputs until its degree becomes large.
Their symmetric product forces a long first row.  Comparing this
trace relation with the free moment complex gives the following
joint restriction on row and label.

\begin{theorem}[A row-and-label strip for positive moment homology]
\label{thm:row-sensitive-moment-homology-strip}
Suppose $\mathbb S_\lambda V\otimes S^\mu$ occurs in
$H_j(\mathcal K_{n,k,D})_D$, where $j\ge1$ and
$\lambda\vdash D$.  Write $a=\lambda_1$.  Then
\begin{equation}\label{eq:row-sensitive-moment-homology-strip}
 n-D+2j\le\mu_1\le
       \left\lfloor\frac{D+a}{2}\right\rfloor-j.
\end{equation}
When $3\le D\le n$, the sharper strip is
\begin{equation}\label{eq:row-sensitive-quadratic-moment-strip}
 n-D+3j\le\mu_1\le
 \min\left\{\left\lfloor\frac{D+a}{2}\right\rfloor-j,\ D-2\right\}.
\end{equation}
More generally, if all moments of orders at most $m\ge1$ form a
regular sequence in $P$, then
\begin{equation}\label{eq:row-sensitive-regular-moment-strip}
 n-D+(m+1)j\le\mu_1\le
 \left\lfloor\frac{D+a-(m+1)(j-1)-2}{2}\right\rfloor.
\end{equation}
\end{theorem}

\begin{proof}
For the general lower bound, remove the linear moments.  Every
remaining exterior generator has weight at least two, and a
polynomial of degree at most $D-2j$ has first label row at least
$n-D+2j$.  In the range $D\le n$, quadratic exactness raises the
minimum remaining weight to three, as in
Lemma~\ref{lem:parabolic-koszul-detection}.

We prove the new upper bound.  For $1\le t\le n$, let $S_t$
act on the last $t$ labels, put $r=n-t$, and write
\[
 A_r=\operatorname{Sym}(V\otimes K^r),\qquad
 A=\operatorname{Sym}V,\qquad
 F=\operatorname{Sym}(A_+),\qquad
 T=(A^{\otimes t})^{S_t}.
\]
The map $F\to T$ sends the formal trace $[b]$ to
$\sum_{i=1}^t b(X_i)$.  By the trace presentation used in
Theorem~\ref{thm:universal-selected-trace-presentation}, its kernel
$I_t$ is generated by the equivariant image of
$\operatorname{Sym}^{t+1}(A_+)$.  Explicitly, the relation on
$b_1,\ldots,b_{t+1}$ is
\[
 \sum_{\pi\in\Pi_{t+1}}
 (-1)^{t+1-|\pi|}
 \prod_{B\in\pi}(|B|-1)!
 \prod_{B\in\pi}\left[\prod_{i\in B}b_i\right].
\]
The internal degree is $\sum_i\deg b_i$; the same formula is
extended multilinearly to arbitrary inputs in $A_+$.

A degree-$g$ generator has at least $2t+2-g$ linear inputs.
Splitting $\operatorname{Sym}^{t+1}(A_+)$ by the number $\ell$
of linear inputs therefore expresses its relevant summands as
$\operatorname{Sym}^{\ell}V\otimes W$, with
$\ell\ge2t+2-g$ and $W$ polynomial.  The
Littlewood--Richardson rule shows that every constituent has first
row at least $2t+2-g$.  Multiplication by any polynomial
representation, or tensoring with one, cannot decrease that first
row.  

First eliminate the global linear moments in this comparison.
They form a regular sequence both in $A_r\otimes F$ and in
$A_r\otimes T=P^{S_t}$.  In the first ring their formal trace
variables are independent.  In the second, splitting the label
representation into the invariant diagonal and the centered
subspace gives a polynomial factor on the global linear moments;
taking $S_t$-invariants preserves this factorization.
Consequently the exact sequence with kernel $A_r\otimes I_t$
remains exact after quotienting by those moments.  Its new kernel
$\overline I_t$ is the image of that ideal, generated by the same
equivariant trace-relation sources.  Their first-row bounds are
therefore unchanged.

The source after elimination is the polynomial algebra
\[
       A_r\otimes\operatorname{Sym}(A_{\ge2}).
\]
In total degrees at most $D$, form its Koszul complex
$\overline{\mathcal L}$ on the remaining global moments
$[X^\alpha]+\sum_{i\le r}X_i^\alpha$, $2\le|\alpha|\le D$.
A polynomial change of the independent formal trace variables
makes this complex exact in positive degrees.  It surjects onto
the corresponding reduced moment complex of $P^{S_t}$, whose
homology is that of $\mathcal K_{n,k,D}^{S_t}$ because the removed
linear sequence is regular.  Its kernel complex has terms
\[
 \mathcal N_i=\overline I_t\otimes
           \bigwedge^i\!\left(\bigoplus_{d=2}^D
                                      \operatorname{Sym}^dV\right).
\]
All maps are equivariant.  Using the full free trace algebra and
then taking the stated degrees gives the same terms: variables
and relations of weight greater than $D$ cannot contribute.

Positive homology of the finite-label complex is detected one
homological degree earlier in this trace-relation kernel.
If a row $\lambda\vdash D$ occurs in $\mathcal N_{j-1}$,
it comes from a trace-relation generator of some degree $g$,
polynomial factors, and $j-1$ exterior generators of weight at
least two.  Hence $g\le D-2j+2$.  By the first-row bound and
Littlewood--Richardson containment,
\[
        a\ge2t+2-g\ge2t+2j-D.
\]
The long exact sequence of
$0\to\mathcal N\to\overline{\mathcal L}\to
\overline{\mathcal K_{n,k,D}^{S_t}}\to0$ identifies positive homology in
degree $j\ge2$ with $H_{j-1}(\mathcal N)$, and injects degree-one
homology into $H_0(\mathcal N)$.  Thus its row-$\lambda$ part
vanishes whenever $D+a<2t+2j$.

Now take $t=\mu_1$.  Branching gives a nonzero $S_t$-fixed
vector in $S^\mu$, and exactness of invariants identifies
$H_j(\mathcal K_{n,k,D})^{S_t}$ with the homology just considered.
Therefore $D+a\ge2\mu_1+2j$, which is the desired upper bound.
This also permits $r=0$: the selected-label polynomial factor is
then $K$.

Combining this estimate with
Lemma~\ref{lem:parabolic-koszul-detection} gives
\eqref{eq:row-sensitive-quadratic-moment-strip}.

If the moments through order $m$ are regular in $P$, their Koszul
complex stays exact after taking $S_t$-invariants, so they are
regular in $P^{S_t}$ as well.  The same comparison can therefore
eliminate all their independent formal traces first.  The
trace-relation sources retain their row bound, while every
remaining exterior generator has weight at least $m+1$.
Thus $g\le D-(m+1)(j-1)$, giving the upper bound in
\eqref{eq:row-sensitive-regular-moment-strip}; removal of the same
regular sequence gives its lower bound by monomial support.
\end{proof}

The \emph{row-$\lambda$ Euler character} is the
$S_n$-representation
\[
 \sum_{j\ge0}(-1)^j
 \operatorname{Hom}_{\operatorname{GL}(V)}
       \bigl(\mathbb S_\lambda V,
              (P\otimes\bigwedge^jU_D)_D\bigr).
\]
Its evaluation at $\sigma\in S_n$ is obtained by projecting the
total-degree-$D$ Koszul product in
\eqref{eq:explicit-stable-character-polynomial} onto the Schur
character $s_\lambda$.  Thus exactness can be tested for one row
partition while other rows in the same multidegree retain homology.



\subsection{Detection of the first row-sensitive trace relations}

At equality in the trace-relation bound, all positive inputs are
linear or quadratic.  This leaves a multiplicity-one source whose
highest vector can be written with bordered minors.  A factorial
evaluation of that vector proves that the first possible relation
survives negative-moment substitution.  It therefore contributes
no positive moment homology.

\begin{lemma}[The boundary trace is detected by selected labels]
\label{lem:boundary-trace-detection}
Let $\lambda\vdash D$, write $a=\lambda_1$, and suppose
\begin{equation}\label{eq:boundary-trace-equality}
                         D+a=2t+2.
\end{equation}
Set $q=t+1$, $r=n-t$, and use the free tail-trace algebra
$F=\operatorname{Sym}(A_+)$, its kernel $I_t$, and the selected-label
polynomial algebra $A_r$ from the proof of
Theorem~\ref{thm:row-sensitive-moment-homology-strip}.
The degree-$D$, row-$\lambda$ part of $A_r\otimes I_t$ has
multiplicity at most one.  It is zero unless
\begin{equation}\label{eq:boundary-even-tail-shape}
       \lambda=(a,2\kappa),\qquad
       |\kappa|=\frac{D-a}{2},\qquad 2\kappa_1\le a.
\end{equation}
If \eqref{eq:boundary-even-tail-shape} holds and
$r\ge\ell(\lambda)$, the negative-moment specialization
\[
 \operatorname{ev}_{-}:A_r\otimes F\longrightarrow A_r,
 \qquad [b]\longmapsto-\sum_{i=1}^r b(X_i),
\]
is nonzero, hence injective, on that single irreducible summand.
Consequently the row-$\lambda$ part of
$H_j(\mathcal K_{n,k,D})^{S_t}_D$ vanishes for every $j\ge1$
if the even-tail condition fails, or if $r\ge\ell(\lambda)$.
A row of length greater than $\dim V$ is absent.
\end{lemma}

\begin{proof}
The trace-relation generators have source
$\operatorname{Sym}^{q}(A_+)$.  A degree-$g$ generator has first
row at least $2q-g$.  Under \eqref{eq:boundary-trace-equality}, a
row-$\lambda$ constituent of total degree $D$ therefore requires
$g=D$, with no positive-degree multiplier.  Moreover its source
must have exactly $a$ linear inputs and
$m=(D-a)/2$ quadratic inputs.  A larger number of linear inputs
would force a larger first row, and any input of degree at least
three would force more linear inputs.  The entire possible source
is consequently
\begin{equation}\label{eq:boundary-trace-source}
     \operatorname{Sym}^{a}V\otimes
          \operatorname{Sym}^{m}(\operatorname{Sym}^{2}V).
\end{equation}

The symmetric-matrix Cauchy identity gives
\[
 \operatorname{Sym}^{m}(\operatorname{Sym}^{2}V)
       =\bigoplus_{\kappa\vdash m}\mathbb S_{2\kappa}V.
\]
Pieri's rule now says that a constituent of
\eqref{eq:boundary-trace-source} with first row exactly $a$ must
be $(a,2\kappa)$, once: its horizontal strip gives
$2\kappa_i\ge\lambda_{i+1}$, while both sides sum to $D-a$.
This proves the asserted multiplicity bound and shape restriction.

It remains to show that this possible summand survives the map
from trace-relation sources to selected polynomials.  Multiplicity
one reduces injectivity to a single nonzero image.  We will choose
a highest vector whose negative-moment image can be evaluated
without cancellation.
Realize
\eqref{eq:boundary-trace-source} using vector generators
$u_1,u_2,\ldots$ of weight one and symmetric-matrix generators
$M_{ij}=M_{ji}$ of weight two.  For $j\ge1$ let
\begin{equation}\label{eq:boundary-bordered-minor}
 p_j(u,M)=
 \det\begin{pmatrix}
  0&u_1&\cdots&u_{j+1}\\
  u_1&M_{11}&\cdots&M_{1,j+1}\\
  \vdots&\vdots&&\vdots\\
  u_{j+1}&M_{j+1,1}&\cdots&M_{j+1,j+1}
 \end{pmatrix}.
\end{equation}
It has $u$-degree two, $M$-degree $j$, and highest weight
$(2^{j+1})$: each raising operator acts by an elementary row
operation together with the corresponding column operation on a
leading principal block.  Set $\kappa_{\ell(\kappa)+1}=0$.
Then
\begin{equation}\label{eq:boundary-highest-vector}
 P_\kappa(u,M)=u_1^{a-2\kappa_1}
       \prod_{j\ge1}p_j(u,M)^{\kappa_j-\kappa_{j+1}}
\end{equation}
is a highest-weight vector of weight $(a,2\kappa)$, with
$u$-degree $a$ and $M$-degree $m$.  For $\kappa=\varnothing$
the formula means $P_\varnothing=u_1^a$.

Introduce auxiliary variables $\xi_1,\ldots,\xi_r$ and the
linear functional
\[
       \mathcal E(\xi_1^{\alpha_1}\cdots\xi_r^{\alpha_r})
                       =\prod_{i=1}^r\alpha_i!.
\]
For the multilinear fundamental trace relation $\Psi_q$, its
negative-moment specialization is exactly
\begin{equation}\label{eq:trace-factorial-moment-identity}
 \operatorname{ev}_{-}\Psi_q(b_1,\ldots,b_q)
  =(-1)^q\mathcal E\!\left[
              \prod_{s=1}^{q}\left(\sum_{i=1}^r
                                  \xi_i b_s(X_i)\right)\right].
\end{equation}
Indeed, all set-partition terms in $\Psi_q$ acquire the same
sign $(-1)^q$.  For a fixed assignment of the inputs to the
selected labels, the sum over set partitions refining each fiber
is the product of the fiber factorials: its blocks are the cycles
of a permutation.  These are precisely the coefficients given
by $\mathcal E$.  Thus the normalization in
\eqref{eq:trace-factorial-moment-identity} is fixed, without any
choice of polarization convention.

Write $\widehat\Psi_q:\operatorname{Sym}^q(A_+)\to I_t$ for
the linear map induced by the multilinear trace relation.  The
actual relation used here is $\widehat\Psi_q(P_\kappa)$.
Applying \eqref{eq:trace-factorial-moment-identity} to it amounts
to substituting
\[
             u=\sum_i\xi_iX_i,\qquad
             M=\sum_i\xi_iX_iX_i^{\mathsf T},
\]
then applying $(-1)^q\mathcal E$.  Since
$r\ge\ell(\lambda)$, specialize $X_i$ to the $i$th coordinate
vector for $1\le i\le\ell(\lambda)$ and set all remaining
$X_i$ equal to zero.  Now $u_i=\xi_i$ and $M$ is diagonal, so
\[
 p_j(u,M)=-\left(\prod_{i=1}^{j+1}\xi_i\right)
                    \left(\sum_{i=1}^{j+1}\xi_i\right).
\]
The resulting $P_\kappa$ is a nonzero polynomial all of whose
nonzero coefficients have sign $(-1)^{\kappa_1}$.
Its image under $\mathcal E$ is therefore a nonzero integer.
This proves nonvanishing over every field of characteristic zero.
The case $\kappa=\varnothing$ gives the same conclusion directly
from $\mathcal E(\xi_1^a)=a!$.

Finally compare the free Koszul complex on
$[X^\alpha]+\sum_{i\le r}X_i^\alpha$ in $A_r\otimes F$
with the moment complex of $P^{S_t}$.  The free complex is exact
in positive degrees; its zeroth homology in the degrees under
consideration is $A_r$.  Its kernel complex has terms
\[
       (A_r\otimes I_t)\otimes\bigwedge^j U_D.
\]
Under \eqref{eq:boundary-trace-equality}, none of its terms with
$j\ge1$ can contain row $\lambda$ in total degree $D$:
an exterior factor would require a trace generator of degree
strictly less than $D$, whose first row is greater than $a$.
All homology in homological degrees at least two is therefore zero
in this row.  The remaining part of the long exact sequence is
\[
 0\longrightarrow H_1(\mathcal K_{n,k,D})^{S_t}_{D,\lambda}
      \longrightarrow
 (A_r\otimes I_t)_{D,\lambda}
             \xrightarrow{\operatorname{ev}_{-}}
                         (A_r)_{D,\lambda}.
\]
The source is zero in the excluded shapes, and in the remaining
case it is a single irreducible on which the displayed map is
nonzero.  This proves the homology assertion.
\end{proof}

\begin{corollary}[The closed upper boundary and row characters]
\label{cor:boundary-row-character-range}
Let $3\le D\le n$, let $\lambda\vdash D$, and put
$a=\lambda_1$.  A joint row-and-label constituent of positive
moment homology satisfies
\begin{equation}\label{eq:boundary-refined-moment-strip}
 n-D+3j\le\mu_1\le
 \min\left\{
       \left\lfloor\frac{D+a}{2}\right\rfloor-j,
       \left\lfloor\frac{D+a-3}{2}\right\rfloor
       \right\}.
\end{equation}
In particular the row-$\lambda$ Koszul Euler character is the
actual row-$\lambda$ character of $(R_{n,k})_D$ whenever
\begin{equation}\label{eq:boundary-row-character-range}
 n\ge\max\left\{D,\;
          D-2+\left\lfloor\frac{D+a-3}{2}\right\rfloor\right\}.
\end{equation}
For a single joint row-and-label component, equality of the Euler
and actual characters already holds, for $3\le D\le n$, if
\begin{equation}\label{eq:boundary-row-label-character-range}
                  \mu_1\ge
                      \left\lfloor\frac{D+a-1}{2}\right\rfloor.
\end{equation}
\end{corollary}

\begin{proof}
The $j$-dependent bounds are
Theorem~\ref{thm:row-sensitive-moment-homology-strip}.
If $D+a$ is odd, they already give the additional upper bound
in \eqref{eq:boundary-refined-moment-strip}.  If $D+a$ is even,
the only additional case to exclude is $j=1$ and
$t=\mu_1=(D+a-2)/2$.  Branching gives an $S_t$-fixed vector.
At that value, Lemma~\ref{lem:boundary-trace-detection} excludes
the row unless $\lambda=(a,2\kappa)$.  In the latter case,
with $m=|\kappa|=(D-a)/2$, the hypothesis $n\ge D$ gives
\[
       r=n-t\ge D-t=m+1\ge\ell(\kappa)+1=\ell(\lambda).
\]
The same lemma therefore excludes the remaining case as well.
This proves the strip.

Write $C=\lfloor(D+a-3)/2\rfloor$.  Under
\eqref{eq:boundary-row-character-range}, positive homology would
require $\mu_1\ge n-D+3\ge C+1$ and $\mu_1\le C$,
a contradiction.  Thus zeroth homology equals the Koszul Euler
character in that row.  The final assertion follows directly
from the upper bound and $C+1=\lfloor(D+a-1)/2\rfloor$.
\end{proof}

For the widest row $a=D$, the strip gives
$n\ge\max\{D,2D-4\}$; for bounded $a$ its threshold is
$\frac32D+O_a(1)$, and balanced three-row partitions have threshold
$\frac53D+O(1)$. These finer rowwise conclusions are independent of
the shorter proof of the uniform theorem.
